%=====================================================================
\chapter{Auctions, Portals and Communities}\label{ch:auctions}
%=====================================================================
\locator{1}{Part I --- E-Commerce and Its Models}

\begin{outcomes}
By the end of this chapter you will be able to:
\begin{tight}
  \item \textbf{Describe} the four common auction formats and \textbf{explain}
        what each one does to a bidder's incentive to bid honestly.
  \item \textbf{Identify} the engineering problems an auction creates that a
        fixed-price shop does not have.
  \item \textbf{Explain} the role of escrow, and \textbf{describe} what a
        marketplace operator must do when two users disagree.
  \item \textbf{Distinguish} a portal, a marketplace and a community, and
        \textbf{state} where each one's revenue comes from.
  \item \textbf{Explain} the two-sided launch problem and \textbf{evaluate}
        strategies for escaping it.
\end{tight}
\end{outcomes}

\begin{prereq}
\chref{models}: the revenue models, especially commission, and the note on
network effects at the end.
\end{prereq}

\begin{keyterms}
auction $\bullet$ English auction $\bullet$ Dutch auction $\bullet$ sealed-bid
auction $\bullet$ Vickrey auction $\bullet$ reserve price $\bullet$ proxy
bidding $\bullet$ sniping $\bullet$ shill bidding $\bullet$ escrow $\bullet$
dispute resolution $\bullet$ reputation system $\bullet$ portal
$\bullet$ marketplace $\bullet$ two-sided market $\bullet$ chicken-and-egg
problem
\end{keyterms}

\section{Why This Chapter Is Here Rather Than at the End}

The published outline lists ``Auctions, Portals, and Communities'' second from
last, after security and deployment. That is where it appears in the book whose
contents page the outline was copied from, and it is the wrong place: these are
not an afterthought to building a shop, they are \emph{alternative shapes} a
business can take, and the decision between them belongs in Part~I alongside
the other models.

So this chapter sits where it belongs. It is also the chapter that explains why
the rest of the book assumes a fixed-price shop: the alternatives are
genuinely harder, and it is worth knowing exactly how.

\section{Letting the Buyer Set the Price}

A fixed-price shop answers ``what is this worth?'' before the customer
arrives. An auction refuses to answer and lets the buyers decide.

That is attractive when the seller genuinely does not know the price --- a
second-hand item, a one-off, something whose value depends on who wants it. It
is unnecessary when the seller does know, which is why most commerce is not
conducted by auction.

\begin{definitionbox}[Four Auction Formats]
\term{English auction} --- ascending, open. The price rises until one bidder
remains. What nearly everyone means by ``auction''.

\smallskip
\term{Dutch auction} --- descending, open. The price starts high and falls
until somebody accepts. Fast, and used for perishable goods where speed matters
more than the last rupee.

\smallskip
\term{Sealed-bid first-price} --- everyone submits one hidden bid; the highest
wins and pays what they bid. Standard for tenders and contracts.

\smallskip
\term{Vickrey auction} (sealed-bid second-price) --- the highest bid wins but
pays the \emph{second}-highest price. Its remarkable property is that bidding
your true valuation is the best strategy whatever anyone else does, so there is
nothing to be gained from strategising.
\end{definitionbox}

\begin{conceptbox}[Why the second-price rule removes the guessing]
Suppose an item is worth Rs.~10{,}000 to you.

\smallskip
In a \emph{first-price} auction, bidding Rs.~10{,}000 guarantees you gain
nothing even if you win --- so you must shade your bid downward, and how far
depends on guessing what others will do. Everyone does this, and everyone
guesses.

\smallskip
In a \emph{second-price} auction, bid Rs.~10{,}000. If the next highest bid is
Rs.~8{,}000, you win and pay Rs.~8{,}000, gaining Rs.~2{,}000. If it is
Rs.~11{,}000, you lose --- and you are glad, because winning would have cost
more than the item was worth to you. Raising your bid above your valuation can
only make you win auctions you did not want; lowering it can only make you lose
auctions you did want. Telling the truth is simply best.

\smallskip
\textbf{Where you have met this already.} The ``proxy bidding'' on consumer
auction sites --- you enter a maximum and the site bids on your behalf up to it
--- is a Vickrey auction wearing an English auction's clothes. The winner pays
one increment above the second-highest maximum, not their own.
\end{conceptbox}

\section{What an Auction Costs You in Engineering}

A fixed-price product is a row that changes rarely. An auction is a small
real-time system, and it brings four problems the rest of this book does not
have to solve.

\rowcolors{2}{white}{lightbg}
\begin{longtable}{@{}L{3.6cm}L{11.4cm}@{}}
\toprule
\rowcolor{primary!15}
\textbf{Problem} & \textbf{Why a fixed-price shop does not have it} \\
\midrule
\endhead
\textbf{Time is part of the state} & A listing must close at an exact moment,
unattended, and settle correctly whether or not anyone is looking. A product
page has no deadline. This needs a scheduler, and a scheduler that did not run
is a silent failure \\
\textbf{Concurrent bids} & Two bids arriving in the same instant must produce
one consistent outcome. This is the classic lost-update problem and it needs a
transaction and a lock, not a quick read-then-write \\
\textbf{Sniping} & A bid placed one second before closing wins without giving
anyone a chance to respond. Fixes exist --- extend the closing time on a late
bid --- and each changes the rules of the auction, which is a design decision
rather than a bug fix \\
\textbf{Shill bidding} & The seller, or a friend, bids to push the price up
with no intention of buying. Detecting it means analysing bidding patterns
across accounts, and it is fraud rather than a feature gap \\
\bottomrule
\end{longtable}

\begin{alertbox}[The lost update, concretely]
Two bidders submit Rs.~5{,}000 at the same instant on an item currently at
Rs.~4{,}500.

\smallskip
A naive implementation reads the current high bid, checks that the new bid
beats it, and writes. Interleaved, both read Rs.~4{,}500, both pass the check,
and both write Rs.~5{,}000 --- the second overwriting the first. One bidder's
money is now owed on an item they did not win, and the audit trail shows a
single bid.

\smallskip
The fix is the one \chref{checkout} uses for stock: do the check and the write
in one atomic statement, inside a transaction, with the row locked. ``Read,
decide, write'' is three steps, and anything can happen between them.
\end{alertbox}

\section{Holding Other People's Money}

An auction site that merely introduces a buyer to a seller has introduced two
strangers and left. Whoever moves first --- sends the money, or sends the goods
--- carries all the risk.

\term{Escrow} is the operator's answer: it holds the buyer's money, tells the
seller to ship, and releases the money when the buyer confirms receipt or a
time limit expires.

\begin{definitionbox}[Escrow and Dispute Resolution]
\term{Escrow} is the holding of a payment by a trusted third party until the
conditions of the transaction are met.

\smallskip
It makes the operator a custodian of money that is not theirs, which in most
jurisdictions --- Pakistan included --- carries regulatory obligations
(\chref{legal}). It also makes the operator the \term{dispute resolution}
authority, because somebody has to decide who is telling the truth when the
buyer says the item never arrived and the seller says it did.

\smallskip
That adjudication cannot be automated away, and it is why marketplace
commissions are higher than payment processing fees: the commission buys
arbitration, not just plumbing.
\end{definitionbox}

\begin{examplebox}[Reputation as the cheap substitute, and what it cannot do]
Escrow is expensive. The cheaper substitute is a \term{reputation system}:
record how every past transaction went and let future counterparties read it.

\smallskip
It works, within limits, and the limits are worth knowing:

\begin{tight}
  \item \textbf{A new seller has no reputation}, so the system gives no
        protection in exactly the case where it is most needed.
  \item \textbf{Reputation can be bought} --- a hundred small honest sales, then
        one large dishonest one. The score is an average; the fraud is a
        maximum.
  \item \textbf{Retaliation distorts it.} If each side rates the other, a buyer
        who reports a problem risks a bad rating in return, so both sides learn
        to stay quiet and the scores drift upward until they mean nothing.
  \item \textbf{It cannot be transferred.} Years of standing on one platform are
        worth nothing on another, which is exactly why the platform likes it.
\end{tight}
\end{examplebox}

\section{Portals, Marketplaces and Communities}

The three are often lumped together and should not be.

A \term{portal} aggregates \emph{information} and sells attention --- a price
comparison site, a news aggregator with shopping links. It carries no stock,
handles no payment, and earns from advertising or affiliate referral. Its
engineering problem is ingesting other people's data and keeping it fresh, and
its commercial weakness is that it owns no transaction and can be cut out the
moment a supplier decides to stop feeding it.

A \term{marketplace} aggregates \emph{sellers} and hosts the transaction. It
earns commission, and it is a shop in which the catalogue is written by
strangers --- which turns catalogue quality, duplicate detection and fraud into
daily operational work rather than a one-time setup.

A \term{community} aggregates \emph{people} who are there for each other rather
than to transact. Commerce, if any, is secondary, and the usual mistake is
assuming an engaged community will convert into a profitable marketplace. Many
do not: the thing people valued was the absence of selling.

\begin{conceptbox}[The two-sided launch problem]
A marketplace is worthless to buyers with no sellers, and worthless to sellers
with no buyers. On day one it has neither, and every new one faces this.

\smallskip
The strategies that actually work are all variations on \emph{refusing to be
two-sided at first}:

\begin{tight}
  \item \textbf{Be one of the sellers yourself} until there are enough others.
        You take on stock you did not want, and you compete with your own
        sellers, and it works.
  \item \textbf{Pick one narrow side and own it completely} --- a single city, a
        single category --- so that the market is small enough to fill.
  \item \textbf{Be useful to one side alone.} Offer sellers a tool worth having
        even with no buyers (stock management, a storefront), then introduce
        buyers once the sellers are already there.
  \item \textbf{Fake the thin side} by listing supply gathered from elsewhere.
        Common, legally risky, and the point at which several well-known
        companies acquired their reputations.
\end{tight}

\smallskip
Notice that none of these is ``build a better marketplace''. The problem is not
a product problem, and a student project that proposes a marketplace without
addressing it has not yet met its hardest part.
\end{conceptbox}

\begin{pitfall}
\begin{tight}
  \item \textbf{Implementing bidding with read-then-write.} Two simultaneous
        bids will eventually produce one lost bid, and it will be
        unreproducible when reported.
  \item \textbf{Closing auctions with a page request.} If closing happens when
        somebody loads the page, an auction nobody looks at never closes. It
        needs a scheduled job, and the job needs monitoring.
  \item \textbf{Treating escrow as a feature.} It is custody of other people's
        money and it is regulated. \chref{legal} is not optional reading if you
        intend to hold funds.
  \item \textbf{Believing a reputation score protects the first transaction.}
        It protects the hundredth.
  \item \textbf{Proposing a marketplace without a plan for the empty side.}
        This is the single most common fatal flaw in student project proposals
        for this course.
  \item \textbf{Assuming a community will monetise.} The engagement you measured
        may exist precisely because nothing was being sold.
\end{tight}
\end{pitfall}

\begin{lab}[Three shapes, and what each one must get right]
No software yet. This is the last of the three analysis laboratories;
\chref{stack} starts building.

\smallskip
Find one live example of each: an \textbf{auction}, a \textbf{marketplace} and
a \textbf{portal}. They may be Pakistani or international, but you must be able
to browse all three.

\smallskip
For each, record:

\begin{tightnum}
  \item \textbf{Who takes the risk} if the goods never arrive --- buyer, seller
        or operator? Find the actual policy page, do not guess. Quote it.
  \item \textbf{Where the money is} between payment and delivery. Is there
        escrow? How long is it held, and what releases it?
  \item \textbf{What the operator earns}, and from which side.
  \item \textbf{One rule that exists because of a specific abuse.} Every mature
        platform has several --- a bidding increment, a minimum feedback score,
        a holding period for new sellers, a limit on listing edits. Name one and
        work out what it is defending against.
\end{tightnum}

\smallskip
Then, for the \textbf{auction} only: find the closing rules. Does a late bid
extend the auction? By how much? Say which auction format that makes it, and
whether it rewards bidding your true valuation.

\smallskip
\textbf{Write up half a page} ending with one sentence: of the three, which
would be hardest to build, and which single problem makes it so?

\smallskip
\textbf{What to notice.} The rules that look arbitrary are nearly always scar
tissue. Each one is a record of something that went wrong, and reading a mature
platform's rules is the cheapest way to learn what will go wrong with yours.
\end{lab}

\begin{checkpoint}
\begin{tightnum}
  \item In a Vickrey auction an item is worth Rs.~15{,}000 to you. Bids of
        Rs.~12{,}000 and Rs.~13{,}500 have been placed by others. What should
        you bid, what do you pay if you win, and why would bidding Rs.~18{,}000
        be a mistake?
  \item An auction site closes listings when a visitor loads the listing page.
        Describe precisely what happens to an auction for an unpopular item,
        and what the seller sees.
  \item A marketplace has plenty of buyers and too few sellers. Give two
        strategies from this chapter that would address it, and name a cost of
        each.
\end{tightnum}
\end{checkpoint}

\begin{chaptersummary}
\begin{tight}
  \item Auctions, portals, marketplaces and communities are alternative shapes
        a business can take, which is why this chapter sits in Part~I rather
        than where the published outline puts it.
  \item The four formats are English, Dutch, sealed-bid first-price and
        Vickrey. Under the Vickrey rule, bidding your true valuation is optimal
        whatever others do --- and consumer proxy bidding is Vickrey in
        disguise.
  \item An auction adds four problems a fixed-price shop does not have: time as
        state, concurrent bids, sniping and shill bidding. The first two are
        engineering; the second two are rule design.
  \item ``Read the high bid, check it, write it'' loses bids under concurrency.
        The check and the write must be one atomic operation.
  \item Escrow makes the operator a custodian of other people's money and the
        arbiter of their disputes. Reputation is the cheap substitute and fails
        exactly on the first transaction and the large one.
  \item Portals aggregate information, marketplaces aggregate sellers,
        communities aggregate people. Only the marketplace owns the
        transaction.
  \item Every two-sided market launches empty, and every strategy that works
        amounts to refusing to be two-sided at the start.
\end{tight}
\end{chaptersummary}

\begin{reviewq}
  \item \textbf{(MCQ)} In a second-price sealed-bid auction, the optimal
        strategy is to bid: (a)~slightly above your valuation; (b)~exactly your
        valuation; (c)~slightly below it; (d)~it depends on the other bidders.
  \item \textbf{(MCQ)} Which problem does a fixed-price shop share with an
        auction site? (a)~Sniping; (b)~shill bidding; (c)~concurrent updates to
        one row; (d)~scheduled closing.
  \item \textbf{(Short)} Explain sniping, and give one rule change that
        discourages it together with the cost of that change.
  \item \textbf{(Short)} Distinguish a portal from a marketplace, and say which
        is more easily cut out of its own business.
  \item \textbf{(Short)} Give three distinct ways a reputation system can be
        gamed or rendered meaningless.
  \item \textbf{(Applied)} Design the database handling for a single bid.
        Give the table columns you would need, the exact check that must happen
        atomically, and state what goes wrong if the check and the write are
        separate statements. You may write it as SQL or as precise prose.
  \item \textbf{(Applied)} A student proposes a marketplace connecting
        Bahawalpur tailors with customers. Identify the two-sided launch
        problem in this specific case, choose one strategy from this chapter,
        and describe concretely what the first eight weeks would consist of.
\end{reviewq}
